%%% EXECUTIVE SUMMARY %%%
\thispagestyle{acknowledgements}

The National Platform for Physical AI Oncology Trials represents a paradigm
shift in how oncology clinical trials are regulated, conducted, and scaled.
Physical AI, defined as the integration of embodied robotic systems with
advanced artificial intelligence for direct patient care in clinical trial
settings, offers a transformative path to treat significantly more patients at
lower cost and higher quality than traditional human-only trial operations. This
platform provides a unified, end-to-end framework spanning regulatory
adaptation, site establishment, patient care, national infrastructure, and
economic analysis. Unlike existing FDA approaches that address AI and robotics
through individual guidance documents issued in isolation, this National
Platform simultaneously adapts three core regulatory standards, provides
complete site establishment documentation, includes validated simulation
evidence at both single-patient and multi-patient scales, defines quantitative
robot readiness standards, and designs national infrastructure for multi-site
coordination.

\vspace{0.3cm}

\noindent\textbf{Regulatory Foundation.}
The platform adapts three internationally recognized regulatory standards to
Physical AI oncology trial contexts. The Adaption: ICH Harmonised Guideline
\cite{ich-adapt} establishes good clinical practice principles for Physical AI
across four major sections covering principles, investigator responsibilities,
sponsor responsibilities, and data governance, building on the ICH E6(R3)
framework \cite{ich-e6r3} recognized across dozens of countries. The Physical
AI Adaption: 21 CFR Part 50 \cite{cfr50-adapt} provides comprehensive human
subjects protection with 17 Physical AI-specific definitions, new informed
consent elements for robotic procedures, and a novel Subpart C dedicated to
Physical AI safety requirements. The Physical AI Adaption: 21 CFR Part 312
\cite{cfr312-adapt} spans ten subparts including a new Subpart J for Physical
AI systems, introduces a Phase 0 simulation validation stage, and establishes
reporting requirements for Physical AI adverse events and cybersecurity
incidents. By building on existing, well-established regulatory frameworks
rather than proposing entirely new legislation, these adaptations add immediate
credibility and practical applicability to Physical AI implementations.

\vspace{0.3cm}

\noindent\textbf{Quantitative Standards.}
Two complementary scoring frameworks ensure responsible robot usage. The
Physical AI Standard Level (PSL), developed within the Clinical Trial Site
Complete Documentation \cite{site-docs}, evaluates site-level regulatory
compliance across three dimensions. The Unification Standard Level (USL)
\cite{usl-standard} evaluates individual robot readiness across four dimensions:
simulation framework switching, AI integration, cross-robot sharing, and
clinical trial collaboration. Scores range from 1.0 to 10.0, with nine robots
evaluated across three categories (cobots, surgical robots, humanoid robots).
The Franka Emika Panda achieved the highest composite score (USL 7.4,
Advanced), while the da Vinci dVRK led surgical robots (USL 7.1, Advanced) with
the highest clinical trial readiness score of any robot evaluated. Together, PSL
and USL create a dual-layer quantitative assessment that replaces subjective
evaluation with reproducible, objective scoring.

\vspace{0.3cm}

\noindent\textbf{Simulation Evidence.}
Two pivotal simulations demonstrate that state-of-the-art AI has the ability to
understand, work with, and scale Physical AI applications beyond human and prior
technology capabilities. A Cancer Patient's Journey, Physical AI
\cite{patient-journey-paper} documents a single-patient journey through all ten
stages of a regulated Physical AI oncology trial, from prescreening through
closeout, for a Stage IIIB non-small cell lung cancer patient over 1,120 days.
The ten-stage pipeline demonstrates end-to-end Physical AI capability with
documented safety outcomes, regulatory compliance across all three adapted
standards, and cost savings at each stage. The 24-hour clinical trial simulation
from the Clinical Trial Site Complete Documentation \cite{site-docs} scaled
operations to 168 patients treated with 29 robots across 15 cancer types,
achieving 99.7 percent uptime with zero patient harm events and 7 adverse events
documented per ICH E6(R3) standards. These simulations establish that more
patients can be treated with more robots and fewer workers, validating the
operational model that underpins the entire platform.

\vspace{0.3cm}

\noindent\textbf{Patient-Centered Design.}
Control has shifted towards the patient's side. Patient Instructions: Physical
AI Trials \cite{patient-instructions-paper} provides accessible documentation
covering interactions with ten robot types. AB 2847 from the site documentation
package establishes legally enforceable patient rights including the right to
refuse robotic procedures, the right to human alternatives, and the right to
real-time information about robot operations. The adapted 21 CFR Part 50
\cite{cfr50-adapt} requires ongoing consent processes beyond the traditional
one-time event, giving patients continuous agency over their participation as
technology may be updated and new robots may be introduced during trials.
Patients benefit from convenience through reduced wait times, lower cost of
treatments through reduced staffing overhead, and higher quality of care through
precision robotics and continuous monitoring.

\vspace{0.3cm}

\noindent\textbf{National Infrastructure.}
Physical AI National MCP Servers \cite{national-mcp-paper} provide
national-scale connectivity through five servers (Governance, Data Flow, Safety,
Interoperability, Analytics) with 23 tools, standardized conformance levels, and
emergency safety modules. Physical AI Federated Learning \cite{fl-paper}
enables privacy-preserving multi-site coordination through five foundational
pillars, a triple AI peer review pipeline using Claude Opus 4.6
\cite{claude-opus}, GPT-5.4 \cite{gpt-5-4}, and Gemini \cite{gemini}, and 82
test files ensuring code trust and safety. Together, these infrastructure
components enable nationwide deployment while maintaining strict patient privacy
under HIPAA \cite{hipaa} and state privacy laws.

\vspace{0.3cm}

\noindent\textbf{Economic Impact.}
Physical AI trial operations offer substantial documented cost advantages. The
single-patient journey documents per-patient cost reductions across all ten
stages, while the 24-hour simulation demonstrates favorable scale economics at
the 168-patient level. Tufts CSDD data \cite{tufts-delay} establishes that each
day of delay in drug development has significant economic consequences. By
compressing trial timelines through 24-hour robotic operations, automated
documentation, and streamlined regulatory processes, Physical AI accelerates the
time to drug approval, translating directly to earlier patient access to
life-saving oncology treatments. The financial case extends beyond cost
reduction to encompass faster approvals, improved data quality, and reduced
regulatory review cycles.

\vspace{0.3cm}

\noindent\textbf{Governance Framework.}
The platform is grounded in the constitutional and statutory framework of the
three branches of U.S. government. Legislative authority through Congressional
budget and oversight mechanisms, executive authority through FDA and HHS
implementation of regulations, and judicial authority through administrative law
review all apply to Physical AI oncology trial governance. Rather than creating
new governance structures, this platform demonstrates how proven mechanisms
ensure transparency, fiscal responsibility, and ethical conduct for Physical AI
trial operations.

\vspace{0.3cm}

This National Platform is more comprehensive than existing FDA approaches
because it: (a) adapts three core regulatory standards simultaneously rather
than issuing piecemeal guidance, (b) provides complete site establishment
documentation across eleven regulatory documents, (c) includes validated
simulation evidence at both single-patient and 168-patient scales, (d) defines
quantitative robot readiness standards through PSL and USL, and (e) designs
national infrastructure for multi-site coordination through MCP servers and
federated learning. All source code is available across three open-source
repositories \cite{main-repo, mcp-repo, fl-repo} under the MIT license,
enabling the pharmaceutical and regulatory industries to evaluate, adopt, and
extend this framework for the advancement of oncology drug development.
